\section*{Devoir maison : La Récursivité en Famille}

On utilise une structure de données pour modéliser une famille, où le nombre total d'oreilles de toutes les personnes (et de leurs descendants) doit être calculé en Python.\\

\hrule

\subsection*{Partie 1 : Structure des Données}

Chaque personne est représentée par un dictionnaire. Ce dictionnaire contient notamment une liste d'\texttt{enfants}, où chaque enfant est lui-même un dictionnaire de personne.

\begin{lstlisting}[language=Python]
famille = {
    'nom': 'Grand-Père',
    'oreilles': 2,
    'enfants': [
        {
            'nom': 'Père',
            'oreilles': 2,
            'enfants': [
                {'nom': 'Moi', 'oreilles': 2, 'enfants': []},
                {'nom': 'Ma Soeur', 'oreilles': 2, 'enfants': []}
            ]
        },
        {
            'nom': 'Tante',
            'oreilles': 2,
            'enfants': [
                {'nom': 'Cousin', 'oreilles': 2, 'enfants': []}
            ]
        }
    ]
}
\end{lstlisting}

\begin{Exercise}[title={Variables et Accès}]

\Question Quel est le \textbf{type} de la variable \texttt{famille} ?\\

\emptybox{1cm}

\Question Quel est le \textbf{type} de \texttt{famille['enfants']} ?\\

\emptybox{1cm}

\Question Donner l'instruction Python permettant de récupérer le \textbf{nom} du premier enfant du 'Grand-Père'.\\

\emptybox{1cm}

\Question Donner l'instruction Python permettant de récupérer le nombre d'\textbf{oreilles} du 'Cousin' (en partant de la variable \texttt{famille}).\\

\emptybox{1cm}

\end{Exercise}

\newpage

\hrule

\subsection*{Partie 2 : Analyse et Implémentation Récursive}

\begin{Exercise}[title={Fonction Récursive}]

\subsubsection*{Question 5 : Compléter la fonction}

L'objectif est de calculer le total des oreilles de toutes les personnes (y compris les descendants). Compléter la fonction Python ci-dessous :

\begin{lstlisting}[language=Python]
def calculer_oreilles_famille(personne):
    # Cas 1 : La personne n'a pas d'enfant.
    if personne['enfants'] == []: # Cas de base
        return ...
    
    # Cas 2 : La personne a des enfants.
    # 1. On compte les oreilles de la personne actuelle.
    total = personne['oreilles']
    # 2. On ajoute les oreilles de tous les descendants.
    for enfant in personne['enfants']:
        total += calculer_oreilles_famille(enfant)
        
    return ...
\end{lstlisting}

\subsubsection*{Question 6 : Rôles}

\Question Quel est le rôle principal du \textbf{Cas 1} (Cas de base) dans la fonction \texttt{calculer\_oreilles\_famille} ?\\

\emptybox{1cm}

\Question Quel est le rôle principal du \textbf{Cas 2} (Cas récursif) ? (Comment ramène-t-il le problème à une instance plus simple ?)\\

\emptybox{1cm}

\subsubsection*{Question 7 : Déroulé}

Considérons cette sous-famille simplifiée :

\begin{lstlisting}[language=Python]
petit_arbre = {'nom': 'A', 'oreilles': 2, 
    'enfants': [
        {'nom': 'B', 'oreilles': 2, 'enfants': []},
        {'nom': 'C', 'oreilles': 2, 'enfants': []}
    ]
}
\end{lstlisting}

\Question Combien de fois au total la fonction \texttt{calculer\_oreilles\_famille} est-elle appelée lors de l'exécution de \texttt{calculer\_oreilles\_famille(petit\_arbre)} ? (Les personnes concernées sont A, B, C.)\\

\emptybox{1cm}

\Question Quel est le résultat final retourné par l'appel \texttt{calculer\_oreilles\_famille(petit\_arbre)} ?\\

\emptybox{1cm}


\end{Exercise}